\chapter{論文執筆の作法}
\label{chap:writing}

本章では，論文を執筆する上での作法をまとめる．
内容・体裁ともにしっかりした論文を書くために，全体の構成から文章表現の細部までを順に説明する．

%% ===========================================================
\section{論文全体の構成}
\label{sec:overall-structure}
論文は，読者が研究の流れを追えるように章を構成する．
実験を伴う研究では，いわゆる IMRAD（Introduction, Methods, Results, and Discussion）型の構成が
よく用いられる．一般的な章構成の例を以下に示す．

\begin{enumerate}
    \item 序論（研究背景・問題提起・目的・論文の構成）
    \item 関連研究（先行研究の整理と本研究の位置づけ）
    \item 提案手法（本研究で提案する手法やシステムの説明）
    \item 評価・実験（実験条件・結果・考察）
    \item 結論（成果のまとめと今後の課題）
\end{enumerate}

上記はあくまで一例であり，研究内容に応じて適切な構成を考えること．
章・節の構成が不適切な論文は，読者にとって大きな負担となる\cite{kisaragi}．

%% ===========================================================
\section{各章の役割と書き方}
\label{sec:each-chapter}

\subsection{序論}
研究の背景と課題を示し，本研究の目的を明確にする．
第\ref{chap:intro}章がその記述例である．

\subsection{関連研究}
本研究に関連する先行研究を整理し，それらとの違いや本研究の新規性を述べる．
既知の技術が何であり，何を新しいアイデアとして提案しているのかを明確にする必要がある．
十分な文献調査と引用は，新規性の主張に欠かせない．

\subsection{提案手法}
提案する手法やシステムを，読者が再現できる程度に具体的に説明する．
概念的・抽象的な説明に終始し，読者が具体的な内容を理解できないものは再考を要する．
図や数式，アルゴリズムを適切に用いて説明するとよい（第\ref{chap:latex}章参照）．

\subsection{評価・実験}
提案手法の有効性を，実験や評価によって客観的に示す．
実験条件を明記し，結果を表やグラフで提示したうえで，その結果が何を意味するのかを考察する．
有効性の主張がない，または極めて貧弱なものは再考を要する．

\subsection{結論}
本研究で行ったことと得られた成果を簡潔にまとめる．
序論で述べた目的に対してどの程度達成できたかを振り返り，今後の課題や展望を述べる．

%% ===========================================================
\section{わかりやすい文章の作法}
\label{sec:writing-style}
研究の新規性・有用性・信頼性が読者に伝わるように記述する\cite{kisaragi}．
読み手にとって読みやすい文章を心がけ，以下の点に注意する．

\begin{itemize}
    \item 「だ・である」調（常体）で統一する．「です・ます」調と混在させない
    \item 一文は短く簡潔にする．極端な口語体や長文の連続は避ける
    \item 主語と述語の対応を明確にする．長い文は主述がねじれやすい
    \item 曖昧な表現（「〜と思われる」「〜的な」）を避け，根拠に基づいて記述する
    \item 指示語（「これ」「それ」）の指す内容が一意に定まるようにする
    \item 説明に飛躍がないか，逆に冗長すぎる部分がないかを確認する
    \item 未定義の用語を減らし，専門用語には適切な説明を付ける
\end{itemize}

%% ===========================================================
\section{句読点と文字の使い分け}
\label{sec:punctuation}
句点には全角の「．」，読点には全角の「，」を用いる（本テンプレートの慣例）．
ただし英文中や数式中では半角の「.」「,」を用いる．「。」「、」は使用しない．

全角文字と半角文字の使い分けは以下の通りである．
\begin{itemize}
    \item 英数字・空白・記号類は半角文字を用いる
    \item カタカナは全角文字を用いる
    \item 括弧は全角の「（）」を基本とし，英文中では半角の「()」を用いる
    \item 引用符は開きと閉じを区別する．開きには\,``\,を，閉じには\,''\,を用いる
\end{itemize}

%% ===========================================================
\section{引用のマナー}
\label{sec:citation-manner}
他者の研究成果や主張を利用する際は，必ず出典を明記する．
出典を示さずに引用することは剽窃であり，固く禁じられている．
本文中での引用には\verb|\cite{}|を用い，引用した文献のみを参考文献リストに列挙する
（記法の詳細は第\ref{sec:reference}節を参照）．

%% ===========================================================
\section{図表の扱い方}
\label{sec:figure-manner}
図表には必ずキャプションを付け，本文中で参照すること．
使用していない図表を掲載してはならない．また，図は原則として自分で作成したものを用いる．
標準テストイメージ等を除き，他者が作成した図表を引用の明示なく用いることは禁止である．
具体的な作成方法は第\ref{sec:figure}節および第\ref{sec:table}節で述べる．

%% ===========================================================
\section{生成AIの利用について}
\label{sec:ai}
ChatGPTやGemini，Claude等の生成AIは，文章の推敲やアイデアの整理など，
執筆を補助するツールとして活用できる．ただし，以下の点に十分注意すること．

\begin{itemize}
    \item 生成AIの出力をそのまま貼り付けてはならない．内容を吟味し，自らの文章として記述する
    \item 生成AIは事実と異なる内容（ハルシネーション）を生成することがある．出力内容は必ず裏付けを取ること
    \item 存在しない参考文献を生成することがあるため，引用情報は必ず原典にあたって確認する
    \item 生成AIの利用方針は指導教員や学科の指示に従うこと
\end{itemize}
内容の責任は著者自身が負うことを忘れてはならない．

%% ===========================================================
\section{推敲と提出前チェック}
\label{sec:proofreading}
論文は必ず提出前に指導教員の確認を受けること．
締切直前にいきなり見せるのではなく，早い段階から計画的に執筆を進め，
フィードバックを受けて改良を重ねる．同級生や先輩に読んでもらうことも有効である．

提出前には以下の点を最終確認する\cite{ipsj-sample}．
\begin{itemize}
    \item \verb|\ThisIsDraft|が削除またはコメントアウトされているか
    \item タイトル・氏名・指導教員名に誤りがないか
    \item 図表番号・参考文献番号が正しく参照されているか（未参照・未定義がないか）
    \item フォントが埋め込まれているか（印刷時の文字化けを防ぐ）
    \item overfull が発生していないか（第\ref{sec:overfull}節参照）
\end{itemize}
